\section{The independent positive-temperature construction}
\label{sec:positive-temperature-alternative}

This appendix constructs near-optimal bounded classical formulas
without using the zero-temperature scalar companion.  Its input is the
positive-temperature Parisi measure at a fixed finite inverse
temperature.  The Parisi diffusion first gives an almost optimal
martingale output.  A finite-step tree calculation and a strong
symmetric-power limit then place that output in the rooted Hilbert
space.  Finally, the bounded approximation and color-exclusion lemmas
of Section~\ref{sec:classical-optimization} put it in the form required
by Theorem~\ref{thm:classical-synthesis}.

Besides providing an independent route to the same synthesis input,
the appendix retains the finite-degree message calculation and the
quantitative fixed-temperature mesh bounds.  Temperature is chosen
first and held fixed in those estimates.  No passage from the
log-cosh terminal datum here to the terminal datum $|x|$ is used.

\subsection{The positive-temperature variational input}

The use of a Parisi diffusion and its gradient martingale to choose
incremental message-passing coefficients follows the classical
construction of Montanari~\cite[Section~3.1, Lemmas~3.2--3.3]{Montanari2019},
with the more general framework of El Alaoui, Montanari, and
Sellke~\cite[Assumption~1 and Theorem~2]{AMS2020}. We use the finite-tree
calculation of El Alaoui, Montanari, and Sellke~\cite{EAMS2023} and
convert its output to the synthesis formulas below.

We record explicitly the spin-glass results that enter the argument.
For an inverse temperature $\beta_0>0$, let $\mu_{\beta_0}$ be the
finite-temperature Parisi measure and put
\[
 \gamma(t)=\beta_0\mu_{\beta_0}([0,t]).
\]
In field coordinates appropriate to the normalization of this paper,
the Parisi PDE and free energy are
\begin{equation}
 \begin{aligned}
 \phi_t(t,x)
   &=-\frac12\bigl(\phi_{xx}(t,x)+\gamma(t)\phi_x(t,x)^2\bigr),\\
 \phi(1,x)&=\frac1{\beta_0}\log\bigl(2\cosh(\beta_0x)\bigr).
 \end{aligned}
 \label{eq:positive-parisi-pde}
\end{equation}
\begin{equation}
 F_{\beta_0}
   =\phi(0,0)-\frac12\int_0^1t\gamma(t)\,\dd t
    =\lim_{n\to\infty}\frac1{n\beta_0}
        \E\log\sum_{\sigma}\exp\bigl(\beta_0h_{n,J}(\sigma)\bigr).
 \label{eq:positive-parisi-value}
\end{equation}
Equation~\eqref{eq:positive-parisi-value} is the positive-temperature
Parisi formula~\cite[Theorem~1.1]{Talagrand2006}, written in the field
coordinates of \eqref{eq:positive-parisi-pde}. The covariance of
$h_{n,J}$ divided by $n$ is $(R_{12}^2-1/n)/2$, where
$R_{12}=n^{-1}\sum_i\sigma_i^1\sigma_i^2$. The $O(1/n)$ term is
allowed by that theorem, so its limiting covariance is $t^2/2$.
The partition function here is the unnormalized sum over configurations.

This route does not need the zero-temperature variational formula,
even for existence of the limit defining $P_*$. Put
$A_n=n^{-1}\E\max_\sigma h_{n,J}(\sigma)$ and
let $F_{n,\beta_0}$ denote the finite-$n$ expression on the right of
\eqref{eq:positive-parisi-value}. The elementary log-sum bounds give
\[
 F_{n,\beta_0}-\frac{\log2}{\beta_0}
 \le A_n\le F_{n,\beta_0}.
\]
For every fixed finite $\beta_0$, the Parisi formula gives convergence
of $F_{n,\beta_0}$. Hence
$\limsup_n A_n-\liminf_n A_n\le(\log2)/\beta_0$.
Letting $\beta_0$ tend to infinity proves convergence of $A_n$ and yields
\begin{equation}
 P_*\le F_{\beta_0}\le P_*+\frac{\log2}{\beta_0}.
 \label{eq:positive-free-energy-bounds}
\end{equation}
A direct zero-temperature variational formula is also known
\cite[Theorem~1]{AC2017}; the argument here needs only the
positive-temperature formula and the preceding squeeze. We use the
following precise finite-temperature facts.  In the conventions of
\cite{Lopatto2026}, they are Theorem~1.1, Lemma~2.2,
Proposition~2.4, and Lemma~2.5, respectively:
\begin{enumerate}
 \item If $\beta_0>1$, there is $q=q_{\beta_0}\in(0,1)$ such that
 $\supp\mu_{\beta_0}=[0,q]$.  The restriction below $q$ has a density
 in $C^\infty([0,q))$; the point $q$ is an atom.
 \item If $X$ solves the stochastic differential equation driven by
 a standard Brownian motion $W$,
 \begin{equation}
  \dd X_t=\gamma(t)\phi_x(t,X_t)\,\dd t+\dd W_t,
  \qquad X_0=0,
  \label{eq:positive-parisi-diffusion}
 \end{equation}
 then $\E\phi_x(t,X_t)^2=t$ for every
 $t\in\supp\mu_{\beta_0}$.
 \item All positive-order spatial derivatives of $\phi$ are bounded
 and jointly continuous on $[0,1]\times\R$.
 The function $\phi$ is even in $x$, $|\phi_x|\le1$, and
 $\phi_{xx}>0$.  These regularity statements also follow from
 \cite[Propositions~1--2 and Appendix Proposition~5, (66)]{AC2015}
 after the same change of variables.
 \item Differentiating the PDE along
 \eqref{eq:positive-parisi-diffusion} gives
 $\dd\phi_x(t,X_t)=\phi_{xx}(t,X_t)\,\dd W_t$.
\end{enumerate}
The self-consistency identity $\E\phi_x(t,X_t)^2=t$ at points
$t\in\supp\mu_{\beta_0}$, including $t=0$, also follows from the
variational argument of Jagannath and Tobasco~\cite[Proposition~1.1 and Corollary~3.10]{JT2017}.
To check the conventions, if $u$ denotes the PDE solution with terminal
condition $\log\cosh x$ and diffusion coefficient $\beta_0^2$ in
\cite{Lopatto2026}, then
$\phi(t,x)=\beta_0^{-1}\bigl(u(t,\beta_0x)+\log2\bigr)$.
In particular, the full-support input used here concerns positive
temperature; no support property of a limiting zero-temperature
minimizer is assumed. The version of \cite{Lopatto2026} used here is
version~3: its Hamiltonian is written with independent Gaussian terms
over all ordered pairs. Combining the two terms for each unordered
pair gives the normalization \eqref{eq:sk-model}; the diagonal terms
form a centered, configuration-independent Gaussian. They have zero
contribution to expected pressure and do not change the Parisi
minimizer.

Write
\begin{equation}
 m(t,x)=\phi_x(t,x),\qquad
 a(t,x)=\phi_{xx}(t,x),\qquad
 b(t,x)=\gamma(t)m(t,x).
 \label{eq:classical-coefficients}
\end{equation}
On every $[0,T]$ with $T<q$, these functions have bounded spatial
derivatives of every order.  They are also Lipschitz in time there:
the density statement makes $\gamma$ smooth on $[0,T]$, and the PDE
expresses the time derivatives of $m$ and $a$ in terms of bounded
spatial derivatives.  In particular, $a$ and $b$ are bounded and
globally Lipschitz in $x$, uniformly on $[0,T]$.
The function $a$ is even in $x$ and $b$ is odd.

\begin{lemma}[A positive-temperature value bound]
\label{lem:positive-temperature-bound}
For every $\beta_0>1$ and $0<T<q$,
\begin{equation}
 \int_0^T\E a(t,X_t)\,\dd t
 \ge P_* -\frac{1/4+\log2}{\beta_0}-(q-T).
 \label{eq:positive-temperature-bound}
\end{equation}
Moreover,
\begin{equation}
 \E a(t,X_t)^2=1\quad(0\le t\le q),
 \qquad 1-q\le\beta_0^{-1}.
 \label{eq:parisi-normalization-and-endpoint}
\end{equation}
\end{lemma}

\begin{proof}
Symmetry gives $m(0,0)=0$.  The martingale identity and It\^o isometry
therefore give
\[
 \E m(t,X_t)^2=\int_0^t\E a(s,X_s)^2\,\dd s.
\]
The left-hand side equals $t$ on $[0,q]$.  Continuity and boundedness
of $a$, together with continuity of $X$, prove the first assertion of
\eqref{eq:parisi-normalization-and-endpoint}, including both endpoints.

Since $\gamma=\beta_0$ on $[q,1]$, the Cole--Hopf formula gives
\begin{equation}
 \phi(q,x)=\frac1{\beta_0}\log\bigl(2\cosh(\beta_0x)\bigr)
                  +\frac{\beta_0(1-q)}2,
 \quad
 m(q,x)=\tanh(\beta_0x),
 \quad
 a(q,x)=\beta_0\operatorname{sech}^2(\beta_0x).
 \label{eq:parisi-endpoint-formula}
\end{equation}
Thus, by Cauchy--Schwarz and the identity just proved,
\[
 1-q=\E\operatorname{sech}^2(\beta_0X_q)
 \le\bigl(\E\operatorname{sech}^4(\beta_0X_q)\bigr)^{1/2}
 =\beta_0^{-1}.
\]

It\^o's formula applied to $X_tm(t,X_t)$ and $\phi(t,X_t)$ yields
\begin{align*}
 \int_0^q\E a(t,X_t)\,\dd t
   &=\E X_qm(q,X_q)-\int_0^qt\gamma(t)\,\dd t,\\
 \E\phi(q,X_q)
   &=\phi(0,0)+\frac12\int_0^qt\gamma(t)\,\dd t.
\end{align*}
The stochastic integrals have mean zero because $a,m,b$ are bounded
and $X$ has finite moments.  The elementary inequality
\[
 x\tanh(\beta_0x)
 \ge\frac1{\beta_0}\log\bigl(2\cosh(\beta_0x)\bigr)
              -\frac{\log2}{\beta_0}
\]
follows, for example, by differentiating the difference for $x\ge0$
and using evenness.  Substitution of
\eqref{eq:parisi-endpoint-formula} consequently gives
\begin{align*}
 \int_0^q\E a(t,X_t)\,\dd t
 &\ge\phi(0,0)-\frac12\int_0^qt\gamma(t)\,\dd t
                -\frac{\beta_0(1-q)}2-\frac{\log2}{\beta_0}\\
 &=F_{\beta_0}-\frac{\beta_0(1-q)^2}{4}
                       -\frac{\log2}{\beta_0}\\
 &\ge P_*-\frac{1/4+\log2}{\beta_0}.
\end{align*}
Finally $0\le\E a(t,X_t)\le(\E a(t,X_t)^2)^{1/2}=1$ on
$[0,q]$, so removing $[T,q]$ costs at most $q-T$.
\end{proof}

\subsection{The finite-iteration tree calculation}

Here we state the finite-step input from El Alaoui, Montanari, and
Sellke with the indices and normalizations needed below.  It is the
general calculation in \cite[Sections~2--4.1]{EAMS2023}, not that
paper's zero-temperature near-optimality theorem.

Fix $\beta_0>1$ and $0<T<q$. Choose $\delta>0$ and an integer
$L\ge1$ with $L\delta\le T$.
Put $t_r=r\delta$.  On the $(D+1)$-regular tree, initialize independent
standard Gaussian labels $g_i$ and set
$u^0_{i\to j}=u_i^0=g_i$, $A^{-1}_{i\to j}=1$.
Each outgoing column multiplies a message by a coefficient from the
preceding step.  This lag will make its newest Gaussian layer
conditionally centered, as verified below.
For $0\le r<L$ use the recursions
\begin{align}
 u^{r+1}_{i\to j}
   &=\frac1{\sqrt D}\sum_{v\in\partial i\setminus\{j\}}
                       A^{r-1}_{v\to i}u^r_{v\to i},
 &u^{r+1}_i
   &=\frac1{\sqrt{D+1}}\sum_{v\in\partial i}
                       A^{r-1}_{v\to i}u^r_{v\to i},
 \label{eq:eams-message-recursion}\\
 x^{r+1}_{i\to j}
   &=x^r_{i\to j}+\delta b(t_r,x^r_{i\to j})
                       +\sqrt\delta\,u^{r+1}_{i\to j},
 &x^0_{i\to j}&=\sqrt\delta\,g_i,
 \label{eq:eams-cavity-euler}\\
 x^{r+1}_i
   &=x^r_i+\delta b(t_r,x^r_i)+\sqrt\delta\,u^{r+1}_i,
 &x^0_i&=\sqrt\delta\,g_i.
 \label{eq:eams-node-euler}
\end{align}
The coefficients, for $0\le r\le L$, are
\begin{align}
 c_{D,r}&=\bigl(\E a(t_r,x^r_{i\to j})^2\bigr)^{-1/2},
 &A^r_{i\to j}&=c_{D,r}a(t_r,x^r_{i\to j}),
 &B^r_i&=c_{D,r}a(t_r,x^r_i),
 \label{eq:eams-finite-normalization}\\
 z_i^L&=\sqrt\delta\sum_{r=1}^LB^{r-1}_iu_i^r.
 \label{eq:eams-output}
\end{align}
The same constant, computed from the directed-edge law, occurs in
$A^r$ and $B^r$.  Its definition is recursive: $x^r_{i\to j}$ has
already been constructed when $c_{D,r}$ is chosen.  Strict positivity
of $a$ ensures every denominator is positive.

The one-step lag in $A^{r-1}u^r$ is important. Let
$\mathcal G_{i\to j}^r$ be the sigma-field of Gaussian labels at
distance at most $r$ from $i$ in the component left after deleting
$\{i,j\}$, including the label at $i$. The variables
$x^{r-1}_{i\to j}$ and $A^{r-1}_{i\to j}$ are
$\mathcal G_{i\to j}^{r-1}$-measurable. The next message is linear
in the Gaussian labels on the new outer layer. More precisely,
\cite[Lemma~3.2]{EAMS2023} gives, for $r\ge1$,
\begin{equation}
 \begin{aligned}
 u^r_{i\to j}\mid\mathcal G_{i\to j}^{r-1}
       &\sim N\bigl(0,(\tau^r_{i\to j})^2\bigr),\\
 (\tau^{r+1}_{i\to j})^2
       &=\frac1D\sum_{v\in\partial i\setminus\{j\}}
       (A^{r-1}_{v\to i})^2(\tau^r_{v\to i})^2,
       \qquad \tau^0=1.
 \end{aligned}
 \label{eq:eams-conditional-innovation}
\end{equation}
The variance recursion also holds at $r=0$. Thus $A^{r-1}u^r$
has conditional mean zero given this earlier cavity history. Its
variance may depend on that history; finite-degree independence of
successive messages is not being asserted. For $r=0$ the sent column
is simply $g_i$.

Let $U_0,U_1,\ldots$ be independent standard Gaussians and define
\begin{align}
 X^\delta_0&=\sqrt\delta\,U_0,
 &X^\delta_{r+1}
  &=X^\delta_r+\delta b(t_r,X^\delta_r)+\sqrt\delta\,U_{r+1},
 \label{eq:eams-limit-euler}\\
 c_{\delta,r}&=\bigl(\E a(t_r,X^\delta_r)^2\bigr)^{-1/2},
 &Z^\delta_L
  &=\sqrt\delta\sum_{r=1}^L
        c_{\delta,r-1}a(t_{r-1},X^\delta_{r-1})U_r.
 \label{eq:eams-limit-output}
\end{align}
At fixed $\beta_0,T,\delta,L$, the update functions and constants
$c_{\delta,r}$ are deterministic. The variables $X^\delta_r$ and
$Z^\delta_L$ are random functions of the displayed Gaussian family.

\begin{lemma}[Finite-step limit]
\label{lem:eams-finite-step}
For the preceding construction, as $D\to\infty$ at fixed $\delta,L$,
\begin{align}
 c_{D,r}&\longrightarrow c_{\delta,r},
 \qquad 0\le r\le L,
 \label{eq:eams-normalizer-limit}\\
 (u^0_i,\ldots,u^L_i)&\Longrightarrow(U_0,\ldots,U_L),
 \qquad z_i^L\Longrightarrow Z^\delta_L,
 \label{eq:eams-distribution-limit}\\
 \lim_{D\to\infty}\frac{D+1}{2\sqrt D}\E z_o^Lz_v^L
   &=\delta\sum_{r=1}^{L-1}
               c_{\delta,r}\E a(t_r,X^\delta_r),
 \qquad o\sim v.
 \label{eq:eams-corrected-value}
\end{align}
The sum in \eqref{eq:eams-corrected-value} is empty when $L=1$.
The distributional statements include convergence of every fixed
moment.  The analogous joint Gaussian limit holds for a directed edge.
Moreover $\E(Z^\delta_L)^2=L\delta$.
\end{lemma}

\begin{proof}
We use the bounded-coefficient CLT and edge calculation in
\cite[Assumptions~2--3, Theorem~3.1, Proposition~2.1, and
Lemmas~4.2--4.3]{EAMS2023}.  The verification below has two parts:
first the degree-dependent normalizers and moment bounds, then the
edge value with its finite-step index.
The functions $a(t_r,\cdot)$ are bounded and Lipschitz, the drift is
Lipschitz, and only finitely many $r$ occur.  Initially
$c_{D,0}=c_{\delta,0}$, and $x^r$ uses normalizers only through
index $r-2$, so the normalization induction is not circular.
Inductively, the Gaussian
limit of the preceding messages and bounded continuity of $a^2$ give
\[
 \E a(t_r,x^r_{i\to j})^2\longrightarrow
 \E a(t_r,X^\delta_r)^2>0.
\]
Thus $c_{D,r}$ converges, and the normalized coefficients and their
Lipschitz constants are bounded uniformly for all sufficiently large
$D$.  This verifies the boundedness and normalization hypotheses used
in the finite-step CLT, including when its coefficients depend on $D$.
One may equivalently run its induction with this triangular sequence:
convergent deterministic coefficients do not affect any induction
step. All fixed moments are uniformly bounded. Indeed, an update
sums the columns $A^{r-1}_{v\to i}u^r_{v\to i}$ from disjoint
child cavity trees. They are independent across children and centered
by \eqref{eq:eams-conditional-innovation}; the bounded multiplier
is measurable before the innovation it multiplies. In this application
their marginal centering also follows from seed parity, since $a$ is
even and $b$ is odd. The moment inequality for sums of independent
centered variables gives, for each integer $k\ge2$,
\[
 \E\left|D^{-1/2}\sum_{j=1}^DY_j\right|^k
 \le C_k\left((\E Y_1^2)^{k/2}
                  +D^{1-k/2}\E|Y_1|^k\right).
\]
Induction over the finite number of updates proves uniform bounds of
every order.  They imply uniform integrability and hence the asserted
moment convergence.  The $x$ and $z$ conclusions follow from their
continuous, polynomial-growth expressions in the messages.

For the energy, put
\[
 R_{D,\ell}=\E\left[
 B_o^{\ell-1}B_v^{\ell-1}
 A_{o\to v}^{\ell-2}A_{v\to o}^{\ell-2}
 u_{o\to v}^{\ell-1}u_{v\to o}^{\ell-1}\right].
\]
Proposition~2.1 of \cite{EAMS2023}, with its degree $k=D+1$, is
the exact identity
\begin{equation}
 \begin{aligned}
 \frac{D+1}{2\sqrt D}\E z_o^Lz_v^L
 &=\delta\sum_{\ell=2}^L
   \E\left[A^{\ell-2}_{v\to o}(u^{\ell-1}_{v\to o})^2
                           B_o^{\ell-1}B_v^{\ell-2}\right]\\
 &\quad+\frac{\delta}{2\sqrt D}\sum_{\ell=1}^L R_{D,\ell}.
 \end{aligned}
 \label{eq:eams-exact-edge-expansion}
\end{equation}
The second line is $O(D^{-1/2})$ by the fixed moment bounds.
To compare node and cavity terms, split off the missing neighbor:
\[
 u_i^r=\sqrt{\frac D{D+1}}\,u_{i\to j}^r
       +\frac{A_{j\to i}^{r-2}u_{j\to i}^{r-1}}{\sqrt{D+1}}
 \qquad (1\le r\le L).
\]
Its squared difference from $u_{i\to j}^r$ is $O(D^{-1})$.
The two Euler recursions and the spatial Lipschitz bounds therefore
give, by induction over this fixed finite number of steps,
\begin{equation}
 \E|x_i^r-x_{i\to j}^r|^2+
 \E|B_i^r-A_{i\to j}^r|^2\le \frac C D.
 \label{eq:eams-node-cavity-bound}
\end{equation}
These are the comparisons in Lemma~4.2 of \cite{EAMS2023}; in
particular, the use of the spatial Lipschitz bound for $a$ is explicit.

Here is the reduction in Lemma~4.3 of that paper in our indices.
Set $s=\ell-1$ and
$H=A_{v\to o}^{s-1}B_v^{s-1}B_o^s$. This product depends only on
the depth-$(s-1)$ cavity labels on the $v$ side and the depth-$s$
cavity labels on the $o$ side. Conditioning on these labels in
\eqref{eq:eams-conditional-innovation} therefore gives
\[
 \E\bigl[H((u_{v\to o}^s)^2-1)\bigr]
 =\E\bigl[H((\tau_{v\to o}^s)^2-1)\bigr].
\]
The variance estimate in \cite[Lemmas~3.2--3.4]{EAMS2023} is uniform
for the present normalizers. Indeed, if $K\ge1$ bounds the
coefficients at the fixed horizon and
$\Delta_r=\E[((\tau^r_{i\to j})^2-1)^2]$, its proof uses only
the variance recursion, independence of child cavities, and the exact
normalization $\E(A^r_{i\to j})^2=1$ to give
\[
 \Delta_{r+1}\le2K^4\Delta_r+\frac{2(K^4+1)}D,
 \qquad \Delta_0=0.
\]
Thus $\Delta_s=O(D^{-1})$, and the preceding conditional expectation
is $O(D^{-1/2})$ by Cauchy--Schwarz. Replacing the two node
coefficients by their cavity counterparts using
\eqref{eq:eams-node-cavity-bound} now yields
\begin{align*}
 \E\bigl[H(u_{v\to o}^s)^2\bigr]
 &=\E\bigl[A_{o\to v}^s(A_{v\to o}^{s-1})^2\bigr]
       +O(D^{-1/2})\\
 &=\E A_{o\to v}^s\,\E\bigl[(A_{v\to o}^{s-1})^2\bigr]
       +O(D^{-1/2})\\
 &=\E B_o^s+O(D^{-1/2}).
\end{align*}
The middle equality uses the exact independence of the two cavity
trees, and the last uses their coefficient normalization and one
more node-to-cavity comparison. This proves the quantitative
replacement without a convergence rate for $c_{D,r}$.
All constants here may depend on $\delta,L,\beta_0,T$, which are fixed.
Taking the degree limit and setting $r=\ell-1$
proves \eqref{eq:eams-corrected-value}.\footnote{Our index convention
is fixed by \eqref{eq:eams-output}.  The finite sums printed in
\cite[(4.5), (4.18)]{EAMS2023} use the shifted endpoints $2,\ldots,L$;
the preceding exact identity and Lemma~4.3 give $1,\ldots,L-1$.
The endpoint difference vanishes in the continuum limit below.
The exact identity we use is Proposition~2.1; the same-time display
in its Lemma~2.8 lacks the factor $1/k$ present in that proposition
and in the lemma's proof.}
Finally, the coefficients multiplying $U_r$ in
\eqref{eq:eams-limit-output} are measurable with respect to
$U_0,\ldots,U_{r-1}$.  Orthogonality of martingale increments and the
definition of $c_{\delta,r-1}$ give
$\E(Z^\delta_L)^2=\delta\sum_{r=1}^L1=L\delta$.
\end{proof}

The degree limit has identified the scalar law $Z^\delta_L$ and the
finite-tree edge limit in \eqref{eq:eams-corrected-value}.  We next
compare these two quantities with the Parisi diffusion at the fixed
temperature.  The symmetric-power argument in the following subsection
will then realize the same finite formula as a rooted vector, so that
its energy can be used in the clipping and synthesis steps.
The normalized Euler and martingale comparison below is related to
\cite[Proposition~5.3]{AMS2020}. We give the estimates for the present
fixed-temperature tree construction.

\begin{lemma}[The continuum limit at fixed temperature]
\label{lem:eams-positive-continuum}
Fix $\beta_0>1$ and $0<T<q$, and let $L=\lfloor T/\delta\rfloor$.
There is a coupling with the diffusion
\eqref{eq:positive-parisi-diffusion} and a constant $C=C(\beta_0,T)$
such that, for all sufficiently small $\delta$,
\begin{align}
 \max_{r\delta\le T}\E|X^\delta_r-X_{r\delta}|^2&\le C\delta,
 &\max_{r\delta\le T}|c_{\delta,r}-1|&\le C\sqrt\delta,
 \label{eq:euler-normalization-bound}\\
 \E|Z^\delta_L-m(T,X_T)|^2&\le C\delta,
 \label{eq:martingale-euler-bound}\\
 \left|\delta\sum_{r=1}^{L-1}c_{\delta,r}\E a(t_r,X^\delta_r)
                    -\int_0^T\E a(t,X_t)\,\dd t\right|
    &\le C\sqrt\delta.
 \label{eq:value-riemann-bound}
\end{align}
\end{lemma}

\begin{proof}
Take $U_0$ independent of $W$ and
$\sqrt\delta\,U_r=W_{r\delta}-W_{(r-1)\delta}$ for $r\ge1$.
Let $\tau_\delta(t)=\delta\lfloor t/\delta\rfloor$ and use the
continuous Euler interpolation
\[
 \overline X^\delta_t
 =\sqrt\delta\,U_0+
       \int_0^tb(\tau_\delta(s),X^\delta_{\lfloor s/\delta\rfloor})\,\dd s
       +W_t.
\]
The drift is bounded and Lipschitz in space and time on $[0,T]$.
Hence
$\E|\overline X^\delta_s-X^\delta_{\lfloor s/\delta\rfloor}|^2
\le C\delta$.
Subtracting the integral equations for $X$ and
$\overline X^\delta$, using Cauchy--Schwarz, and then Gronwall gives
\[
 \E\sup_{s\le T}|\overline X^\delta_s-X_s|^2\le C\delta.
\]
The initial contribution is $\E|\sqrt\delta U_0|^2=\delta$.
This proves the first bound of \eqref{eq:euler-normalization-bound}.

Because $a$ is bounded and spatially Lipschitz, and
$\E a(t_r,X_{t_r})^2=1$,
\[
 \max_{r\delta\le T}
 \left|\E a(t_r,X^\delta_r)^2-1\right|\le C\sqrt\delta.
\]
For small $\delta$ all these expectations lie in $[1/2,3/2]$.
The map $x\mapsto x^{-1/2}$ is Lipschitz on that interval, proving
the second bound of \eqref{eq:euler-normalization-bound}.

Equation~\eqref{eq:eams-limit-output} is the stochastic integral
\[
 Z^\delta_L=\int_0^{L\delta}
 c_{\delta,\lfloor t/\delta\rfloor}
 a(\tau_\delta(t),X^\delta_{\lfloor t/\delta\rfloor})\,\dd W_t.
\]
The integrand differs from $a(t,X_t)$ in squared mean by at most
$C\delta$, uniformly in $t\le L\delta$.  This follows from
\eqref{eq:euler-normalization-bound}, the spatial and temporal
Lipschitz bounds for $a$, and
$\E|X_t-X_{\tau_\delta(t)}|^2\le C\delta$.
It\^o isometry, including the last interval of length at most
$\delta$, gives \eqref{eq:martingale-euler-bound}, since
$m(T,X_T)=\int_0^Ta(t,X_t)\,\dd W_t$.
The same estimates in first mean compare the sum in
\eqref{eq:value-riemann-bound} with its Riemann integral.
The omitted first and last intervals have total length at most
$3\delta$ and bounded integrands.  This proves the final estimate.
\end{proof}

\subsection{Strong convergence of classical descendant formulas}

We next place this classical construction in the Hilbert space used
by Proposition~\ref{prop:energy}.  This step requires norm convergence,
not only a CLT in distribution.

We use the Gaussian lift $\mathsf G$ of
Section~\ref{subsec:classical-gaussian-sector} and the finite descendant
formulas of Definition~\ref{def:finite-descendant-formula}.
In the directed
EAMS calculation the sent columns are $A^{r-1}u^r$, with $u^0=g$
for the first propagation.  The node output uses these same columns.
Replacing $D+1$ root children by $D$, or the final normalization
$\sqrt{D+1}$ by $\sqrt D$, does not change the limiting formula.

\begin{lemma}[The symmetric-power realization]
\label{lem:classical-strong-limit}
Consider a fixed finite descendant formula whose pointwise functions
are smooth with all derivatives of polynomial growth.  Assume its
leaf seeds have moments of every order and every aggregated child
column has mean zero under its finite-degree child law for each $D$.
Its finite-degree root vector converges in
norm, under the symmetric occupation embeddings, to the formula
obtained by replacing every normalized sum by $\mathsf G$.
The same conclusion holds for a convergent finite collection of
deterministic coefficients in the formula.  The finite-degree root
vectors have a uniformly bounded first child-occupation moment.
Consequently, for centered equivariant outputs $f_D$ and their limit
$f$,
\begin{equation}
 \lim_{D\to\infty}\frac{D+1}{2\sqrt D}\E f_D(o)f_D(v)
       =\cE(f),\qquad o\sim v.
 \label{eq:classical-energy-transfer}
\end{equation}
\end{lemma}

\begin{proof}
The proof first treats characteristic vectors, then arbitrary
pointwise maps, and finally the occupation bound needed for energy.
We use induction on formula depth.  All intermediate
moments are uniformly bounded: normalized centered sums satisfy the
moment bound used in Lemma~\ref{lem:eams-finite-step}, and
polynomial-growth functions preserve the existence of sufficiently
high uniform moments.

Suppose the real centered child columns
$f_{D,1},\ldots,f_{D,k}$ already converge in norm to
$f_1,\ldots,f_k$ in the common child Hilbert space.  Fix
$t\in\R^k$ and put $F_D=\sum_jt_jf_{D,j}$,
$F=\sum_jt_jf_j$. The limits are centered because the embeddings
preserve the vacuum and the convergence is in norm.
A joint characteristic vector for the parent sums
is the symmetric tensor power of
\[
 h_D=\exp(\ii F_D/\sqrt D)\vac
          =\alpha_D\vac+\zeta_D,
 \qquad \zeta_D\perp\vac.
\]
Taylor expansion in $L^2$, justified by the uniform fourth moment,
and centering give
\[
 \alpha_D=1-\frac{\norm{F_D}^2}{2D}+O(D^{-3/2}),
 \qquad \sqrt D\,\zeta_D\longrightarrow\ii F.
\]
The $j$-particle component of the embedded tensor power is
\[
 \binom Dj^{1/2}\alpha_D^{D-j}\zeta_D^{\otimes j}
 \longrightarrow
 e^{-\norm F^2/2}\frac{(\ii F)^{\otimes j}}{\sqrt{j!}}.
\]
These are the components of $e^{\ii q(F)}\vac$.
The expected occupation of the left-hand vector is
$D\norm{\zeta_D}^2=O(1)$, so its occupation tails are uniformly
small.  Thus convergence of each fixed grade implies norm convergence
of the complete vector.  A root with $M=D+1$ children has the same
proof: use $M$ factors and the normalization $M^{-1/2}$, with the
same convergent child columns.  Alternatively, retaining $D^{-1/2}$
at that root gives the same limit because $M/D\to1$.  The occupation
bound in either convention is $M\norm{\zeta_D}^2=O(1)$.

Finite linear combinations of these exponentials, multiplied by
bounded measurable functions of the independent root seed, are dense
in the limiting $L^2$ space.  Indeed the root seed and aggregate
fields have a product law.  Products of dense families in the two
factors are dense, and exponentials are dense in the field factor:
an orthogonal $L^2$ vector would give a finite signed measure with
zero Fourier transform and would therefore vanish.
There is a useful precision concerning measurable seed factors.
Let $S$ denote the root seed, $Y_D$ the finite list of aggregate
fields, and $Y$ its limit.  The law of $S$ is unchanged with $D$ and
is independent of $Y_D$ and $Y$.  For
$P(s,y)=\sum_jg_j(s)e^{\ii t_j\cdot y}$ with bounded measurable
$g_j$, and a bounded target $F(s,y)$ continuous in $y$ for every
fixed $s$, the function $|F(s,y)-P(s,y)|^2$ is bounded continuous in
$y$ for each fixed $s$.  Weak convergence of $Y_D$ therefore gives
convergence of its field integral pointwise in $s$.  The bound is
uniform in $s$, so dominated convergence in the unchanged root-seed
law proves
\[
 \E|F(S,Y_D)-P(S,Y_D)|^2
       \longrightarrow\E|F(S,Y)-P(S,Y)|^2.
\]
Thus measurability of $g_j$ does not require a discontinuous test
function in an appeal to joint weak convergence.  For a
polynomial-growth target, first multiply by smooth cutoffs outside
large boxes in the seed and field variables.  The cutoffs preserve
continuity in the field variables, and uniform higher moments make
the discarded squared norms uniformly small.  This proves the induction
step.  It also proves
continuity in deterministic coefficients, by applying the same
argument to a convergent sequence of them.

For completeness, the occupation bound is not a consequence of the
CLT alone.  At a parent, let $P_j$ average over the $j$th child, and
write its output as $H_D(s,S_{D,1},\ldots,S_{D,k})$, where
$S_{D,l}=D^{-1/2}\sum_j f_{D,l}^{(j)}$.
The occupation operator in the symmetric embedding is
$N_D=\sum_{j=1}^D(I-P_j)$.  With an independent replacement of child
$j$, conditional variance and the mean-value theorem give
\[
 \norm{(I-P_j)H_D}^2
 \le\E|H_D-H_D^{(j)}|^2\le\frac{C}{D}.
\]
Here the derivative of $H_D$ has polynomial growth and the sums,
the replaced child columns, and their products have the uniform
moments established above; H\"older's inequality therefore gives a
constant independent of $j,D$.  Summing proves
$\braket{H_D}{N_DH_D}\le C$.  If needed, the same argument with
distinct child replacements and mixed derivatives bounds the
factorial moments of $N_D$ of every fixed order.

The energy identity is now Proposition~\ref{prop:energy} applied to
these classical vectors.  More explicitly, with $M=D+1$ root
children, the projection onto an excitation in a distinguished child
has norm at most $C/\sqrt M$ by symmetry and the occupation estimate.
The term with excitations across the distinguished edge on both sides
is therefore $O(M^{-1})$.  The two terms with only one crossing
select a single child; after multiplication by $\sqrt M$ they
converge to $\braket{f}{\cL f}$ and its conjugate.  The term with
neither crossing is zero by centering.  The factor
$M/(2\sqrt D)$ in \eqref{eq:classical-energy-transfer} thus gives
$\re\braket{f}{\cL f}$, since $M/D\to1$.
\end{proof}

Apply this lemma to the EAMS formulas at fixed $\delta,L$.
All sent columns are odd under simultaneous reversal of the seed
labels, since $a$ is even and $b$ is odd. They are therefore centered
under each finite-degree child law, as required by the lemma.
The smooth polynomial-growth hypotheses hold because a transmitted
column is a bounded smooth coefficient times a message
variable.  Equation~\eqref{eq:eams-normalizer-limit} and the
coefficient-continuity assertion allow us to replace $c_{D,r}$ by the
fixed limiting constants $c_{\delta,r}$ in the limiting formula.
Let its output vector be $z^\delta$.  We obtain
\begin{equation}
 \cE(z^\delta)=\delta\sum_{r=1}^{L-1}
                    c_{\delta,r}\E a(t_r,X^\delta_r),
 \quad
 \norm{z^\delta}^2=L\delta,
 \quad
 z^\delta\ \hbox{has law }Z^\delta_L.
 \label{eq:classical-eams-vector}
\end{equation}
All coefficients in this vector are independent of degree.

\subsection{Clipping and finite bounded formulas}

Let $\pi(x)=\max(-1,\min(x,1))$.  Since $|m(T,X_T)|\le1$,
the law in \eqref{eq:classical-eams-vector} and the metric projection
property give
\[
 \norm{z^\delta-\pi(z^\delta)}^2
 =\E|Z^\delta_L-\pi(Z^\delta_L)|^2
 \le\E|Z^\delta_L-m(T,X_T)|^2.
\]
Only the marginal law of $Z^\delta_L$ enters the equality. The
diffusion coupling bounds this scalar clipping error; the energy of
$z^\delta$ has already been identified separately in the rooted space
by \eqref{eq:classical-eams-vector}.
Consequently, by Lemma~\ref{lem:eams-positive-continuum},
$\norm{z^\delta-\pi(z^\delta)}\le C\sqrt\delta$.
We apply the shared energy-continuity estimate
\eqref{eq:energy-l2-continuity} with $f=z^\delta$ and
$g=\pi(z^\delta)$; both norms are at most one.  Combining
\eqref{eq:positive-temperature-bound},
\eqref{eq:value-riemann-bound}, and
\eqref{eq:classical-eams-vector} proves the following: for every
$\varepsilon>0$, first choosing a sufficiently large finite
$\beta_0$, then $T<q$ sufficiently close to $q$, and finally a
sufficiently small positive $\delta$, gives a centered odd classical
vector
\begin{equation}
 m^*=\pi(z^\delta),\qquad |m^*|\le1,
 \qquad \cE(m^*)\ge P_*-\varepsilon.
 \label{eq:clipped-classical-vector}
\end{equation}
At this point the depth $L$ and the entire descendant calculation are
finite and fixed.

The next proposition completes this route using only the common
bounded-approximation and exclusion lemmas.  It gives the same
conclusion as Proposition~\ref{prop:classical-near-optimum}, with an
independent proof from the clipped vector just constructed.

\begin{proposition}[The independent positive-temperature construction]
\label{prop:classical-near-optimum-positive-temperature}
For every $\varepsilon>0$ and every $\eta_H>0$, there are a finite
color partition with helper probability $0<p_H<\eta_H$, a finite
number of helper types, and finite bounded odd angle formulas
$\theta_c$ on the nonhelper colors such that the following hold.
Each transmitted column is bounded and odd, every pointwise map
may be chosen as a finite trigonometric sum in its odd inputs, and
every aggregation excludes its current color, all path-ancestor
colors, and all helpers.  Conditional on the finite labels, all
continuous derivatives of the pointwise maps are bounded.  The
output
\begin{equation}
 m^{\mathrm{col}}
  =\sum_{c=1}^M\1_{\{C=c\}}\sin(2\theta_c)
 \label{eq:positive-final-classical-formula}
\end{equation}
is centered, satisfies $|m^{\mathrm{col}}|\le1$, vanishes on helper
roots, and obeys
\begin{equation}
 \cE(m^{\mathrm{col}})\ge P_*-\varepsilon.
 \label{eq:positive-final-classical-value}
\end{equation}
After the uncolored finite formula has been fixed, the total helper
probability may be taken arbitrarily small, with all individual
helper probabilities positive.  All depths, coefficients, label laws,
and functions in the conclusion are fixed independently of degree,
system size, and disorder.
\end{proposition}

\begin{proof}
Use \eqref{eq:clipped-classical-vector} with error
$\varepsilon/3$.  Its underlying formula $z^\delta$ has odd
transmitted columns and jointly odd smooth local maps with
polynomial-growth derivatives, as checked after
Lemma~\ref{lem:classical-strong-limit}.  It therefore meets the
hypotheses of Lemma~\ref{lem:classical-trigonometric-density}.
Apply that lemma and then Lemma~\ref{lem:classical-color-exclusion}
with the choices in the proof of
Proposition~\ref{prop:classical-near-optimum}: the first output error
is less than $\varepsilon/12$ in norm, the second less than
$\varepsilon/6$, and one positive-probability helper is reserved per
propagation depth with total probability below $\eta_H$.
The two further energy losses are less than $\varepsilon/6$ and
$\varepsilon/3$, respectively, by
\eqref{eq:energy-l2-continuity}.  Their sum with the initial loss is
less than $\varepsilon$.  Those two shared lemmas also give the
boundedness, trigonometric form, exclusions, and parity asserted here.
The entire formula and all helper probabilities are fixed after the
degree limit of the finite calculation.
\end{proof}
